\chapter{概述与泄露始末}

\section{Claude Code 是什么}

Claude Code 是 Anthropic 公司开发的官方 CLI（命令行界面）编程助手。它允许开发者直接在终端中与 Claude 大语言模型交互，执行软件工程任务——包括编辑文件、运行命令、搜索代码库、管理 Git 工作流等。

与 ChatGPT 等网页界面不同，Claude Code 的定位是\textbf{嵌入开发者工作流}的原生工具。它不是一个简单的聊天机器人，而是一个拥有完整工具链的 AI 编程代理（Agent）。它可以：

\begin{itemize}[nosep]
  \item 读写文件系统中的任意文件
  \item 执行 Shell 命令（Bash/PowerShell）
  \item 使用 ripgrep 搜索代码
  \item 创建 Git 提交和 Pull Request
  \item 生成子代理（Sub-agent）并行处理任务
  \item 通过 MCP（Model Context Protocol）连接外部工具
  \item 在 VS Code 和 JetBrains 等 IDE 中作为后端运行
\end{itemize}

\section{泄露事件}

\subsection{时间线}

2026 年 3 月 31 日，安全研究员 \href{https://x.com/Fried_rice}{Chaofan Shou (@Fried_rice)} 在 X（原 Twitter）上公开了一个发现：

> "Claude code source code has been leaked via a map file in their npm registry!"

\subsection{泄露机制}

泄露的技术细节如下：

\begin{enumerate}[nosep]
  \item \textbf{Source Map 文件}：Anthropic 在将 Claude Code 发布到 npm 注册表时，发布包中包含了一个 \code{.map}（Source Map）文件。
  \item \textbf{Source Map 中的引用}：该 Source Map 文件中包含了对完整、未混淆的 TypeScript 源码的引用路径。
  \item \textbf{R2 存储桶}：这些源码文件存储在 Anthropic 的 Cloudflare R2 存储桶中，且可以公开下载。
  \item \textbf{完整源码可获取}：通过解析 Source Map 中的路径，可以下载整个 \code{src/} 目录的完整源码。
\end{enumerate}

\begin{lstlisting}[language=TypeScript,breaklines=true]
npm package (.tgz)
  └── dist/
       └── bundle.js.map  ← Source Map 文件
            └── sourcesContent 引用 → R2 存储桶中的 TypeScript 源文件
\end{lstlisting}

\subsection{影响范围}

泄露的源码包含：
\begin{itemize}[nosep]
  \item 完整的 \code{src/} 目录，约 \textbf{1,900 个文件}
  \item 总计 \textbf{512,000+ 行} TypeScript 代码
  \item 所有核心模块的完整实现
\end{itemize}

\textbf{未包含的部分}：
\begin{itemize}[nosep]
  \item \code{@ant/*} 开头的 Anthropic 内部专用包（如 \code{@ant/computer-use-mcp}、\code{@ant/claude-for-chrome-mcp}）
  \item 部分构建脚本和 CI/CD 配置
  \item 测试文件（部分）
\end{itemize}

在公开仓库中，缺失的内部包被替换为导出空操作（no-op）的 stub 模块。

\section{项目基本信息}

\begin{longtable}{ll}
\toprule
属性 & 值 \\
\midrule
语言 & TypeScript（严格模式） \\
运行时 & \href{https://bun.sh}{Bun} v1.3+ \\
终端 UI 框架 & React + \href{https://github.com/vadimdemedes/ink}{Ink} \\
CLI 解析 & \href{https://github.com/tj/commander.js}{Commander.js}（extra-typings） \\
模式验证 & \href{https://zod.dev}{Zod} v4 \\
代码搜索 & \href{https://github.com/BurntSushi/ripgrep}{ripgrep} \\
协议支持 & MCP SDK、LSP \\
API 客户端 & Anthropic SDK \\
遥测 & OpenTelemetry + gRPC \\
功能开关 & GrowthBook \\
认证 & OAuth 2.0、JWT、macOS Keychain \\
\bottomrule
\end{longtable}

\section{架构概览}

Claude Code 的整体架构可以用一张图来概括：

\begin{lstlisting}[language=TypeScript,breaklines=true]
用户输入 (终端)
    │
    ▼
┌─────────────────────────────────────────────┐
│  main.tsx (Commander.js CLI 解析)            │
│  ├── 启动优化：并行预取 MDM / Keychain       │
│  └── React/Ink 渲染器初始化                  │
└─────────────────┬───────────────────────────┘
                  │
    ▼─────────────┼──────────────▼
┌────────────┐    │    ┌────────────────┐
│  REPL 模式  │    │    │  Print 模式     │
│ (交互式)    │    │    │ (单次执行 -p)   │
└─────┬──────┘    │    └───────┬────────┘
      │           │            │
      ▼───────────┴────────────▼
┌─────────────────────────────────────────────┐
│         QueryEngine (查询引擎)               │
│  ├── 构建 System Prompt                      │
│  ├── 管理消息历史                             │
│  ├── 调用 Anthropic API                      │
│  ├── 处理 Tool Use 循环                      │
│  └── Token 计数 & 成本追踪                   │
└─────────────────┬───────────────────────────┘
                  │
    ▼─────────────┼──────────────▼
┌────────────┐    │    ┌────────────────┐
│  工具系统   │    │    │  命令系统       │
│ (40+ 工具)  │    │    │ (50+ 斜杠命令) │
└─────┬──────┘    │    └───────┬────────┘
      │           │            │
      ▼───────────┴────────────▼
┌─────────────────────────────────────────────┐
│              服务层 (Services)                │
│  ├── Anthropic API 客户端                    │
│  ├── MCP 服务器连接                          │
│  ├── OAuth 认证                              │
│  ├── LSP 语言服务                            │
│  ├── 上下文压缩 (Compact)                    │
│  ├── 分析 / 遥测                             │
│  └── 插件 / 技能加载                         │
└─────────────────────────────────────────────┘
\end{lstlisting}

\section{本书的分析方法}

本书采用\textbf{自顶向下}与\textbf{自底向上}相结合的方法：

\begin{enumerate}[nosep]
  \item \textbf{自顶向下}：从入口文件 \code{main.tsx} 开始，沿着启动流程和请求处理链路，逐步深入到每个子系统。
  \item \textbf{自底向上}：对关键的类型定义（如 \code{Tool.ts} 中的接口）和工具实现（如 \code{BashTool}）进行独立的深度分析。
  \item \textbf{横向对比}：将 Claude Code 的设计决策与业界常见做法（如 LangChain、AutoGPT 等）进行对比，分析其设计取舍。
\end{enumerate}

每章结尾都包含\textbf{设计启示}部分，总结该模块中值得借鉴的设计模式和工程实践。

\bigskip\noindent\rule{\textwidth}{0.4pt}\bigskip

\textbf{下一章}：\href{./chapter-02-tech-stack.md}{技术栈与项目结构}
